%! TeX root = ./main.tex
% vim: set spelllang=fr:
\chapter*{Introduction}
\addcontentsline{toc}{chapter}{Introduction}
\markboth{\textsc{introduction}}{}

% Phylogénétique, introduite par l’exemple de l’évolution des gènes
% Exemples de motivations
Les gènes sont les véhicules de l’hérédité chez les êtres vivants.
Ils contribuent à définir leurs caractères et peuvent être transmis d’une génération à la suivante.
Au fil du temps et au gré de divers événements, ils peuvent évoluer en différents allèles et être disséminés à travers plusieurs espèces.
La phylogénétique vise à retracer l’histoire de l’évolution d’unités telles que les gènes, les cellules ou les espèces.
Ces histoires sont des sources précieuses d’information pouvant permettre d’élucider les fonctions de gènes~\parencite{Lynch2001ProbabilityPreservationNewly}, de caractériser l’évolution de tumeurs~\parencite{Quinn2021SingleCellLineages} ou celle de pathogènes~\parencite{Morel2020PhylogeneticAnalysisSars}.

% Arbres phylogénétiques et événements
Un \emph{arbre phylogénétique} représente une hypothèse d’histoire évolutive.
Comme un arbre généalogique, un arbre phylogénétique est composé d’une suite d’embranchements et de branches le long desquelles s’écoule le temps jusqu’à rejoindre les représentants contemporains.
Chaque embranchement de l’arbre correspond à la divergence putative d’un ancêtre en plusieurs descendants libres d’évoluer indépendamment.
Les causes possibles de telles divergences diffèrent selon les types d’unités étudiés; par exemple, les gènes se dupliquent, les cellules se divisent et les espèces subissent des spéciations.

% Problème d’évolution hiérarchique
L’évolution d’un gène ne dépend pas seulement des événements qui l’affectent directement, mais également de l’évolution de son espèce hôte.
La spéciation d’une espèce cause en effet l’émergence d’exemplaires indépendants de ses gènes, un dans chaque nouvelle espèce.
% Cette relation est asymétrique --- la seule duplication d’un gène ne pouvant pas causer une spéciation --- et forme ainsi une hiérarchie évolutive entre un ensemble d’\emph{unités dépendantes} et leurs hôtes.
Cette relation forme une hiérarchie évolutive entre un ensemble d’\emph{unités dépendantes} et leurs hôtes.
Des hiérarchies similaires se retrouvent à d’autres échelles, par exemple entre les domaines protéiques et leur protéine, entre certains parasites et leur hôte, ou entre les espèces et leur habitat\footnote{Mais aussi en dehors de la biologie, comme entre les fragments de code et leur programme. Un tel fragment peut être copié-collé à l’intérieur d’un programme, et une fourche (\emph{fork}) peut avoir lieu dans le programme hôte. Dans les deux cas, des copies indépendantes de fragments apparaissent et continuent d’évoluer indépendamment par la suite.}.

% Réconciliation
Pour étudier ces hiérarchies évolutives, \textcite{Goodman1979FittingGeneLineage} ont introduit la méthode dite de \emph{réconciliation}.
En s’appuyant sur les différences et similarités topologiques entre la phylogénie d’une famille de gènes et celle de ses espèces hôtes, la réconciliation permet d’établir à laquelle des deux se rapporte chaque divergence.
Le résultat est un arbre réconcilié, synthèse des deux arbres, dans lequel chaque embranchement de l’arbre de gènes est annoté par le type d’événement évolutif l’ayant engendré.

% Problème d’évolution segmentale
Bien qu’elle permette l’étude de hiérarchies évolutives entre des unités dépendantes et leurs hôtes, la réconciliation suppose que les unités étudiées évoluent indépendamment les unes des autres.
Cette hypothèse est invalidée en particulier lorsque plusieurs unités sont organisées linéairement dans le même hôte, comme des gènes le long d’un génome, et affectées par des événements à grande échelle.
Différentes méthodes ont été proposées pour étendre le paradigme à une réconciliation multiple, prenant en considération l’évolution conjointe de plusieurs unités.

% Projet de recherche
L’objectif général du projet de recherche présenté dans ce document consiste à développer ces méthodes de réconciliation multiple.
Il s’agit notamment de rechercher un équilibre entre puissance de modélisation et complexité algorithmique.
L’application de ces méthodes à un cas biologique est également considérée pour les systèmes synténiques Crispr-Cas.

% Plan du rapport
Les deux premiers chapitres de ce rapport sont consacrés à une revue de la littérature pertinente.
Dans le premier chapitre, une synthèse de l’état de l’art des méthodes de réconciliation est présentée.
Le deuxième chapitre expose les généralisations existantes de la réconciliation pour l’étude de multiples unités dépendantes organisées linéairement.
Enfin, le dernier chapitre décrit les axes principaux du projet de recherche, en faisant état du travail déjà accompli et en dressant un échéancier prévisionnel.

%%%%%%%%%%

%Le concept de gène a été postulé dès~1865 par Mendel~\cite{Bateson1902MendelsPrinciplesHeredity}, mais sa compréhension s’est considérablement développée au cours du siècle suivant, notamment grâce aux découvertes en biologie moléculaire.
%La seconde moitié du XX\textsuperscript{e} siècle a en effet vu l’élucidation de la structure et des fonctions de l’acide désoxyribonucléique (ADN).
%Cette longue molécule, présente dans le noyau de chaque cellule, est porteuse de l’information du génome, qui se manifeste dans la séquence de nucléotides qui la composent.
%Ces nucléotides se déclinent en quatre variétés qui diffèrent par leur base azotée : l’adénine~(A), la cytosine~(C), la guanine~(G) et la thymine~(T).
%Chaque gène est caractérisé par un fragment unique de cette séquence, qui peut être vue comme une chaîne de caractères choisis parmi quatre lettres : A, C, G et T.
%Avec le temps, cette séquence est soumise à des mutations à différentes échelles, qui sont filtrées par la sélection naturelle, et contribuent à l’évolution des gènes en plusieurs versions.

%Pour ce faire, il peut être pertinent d’observer les similitudes et différences dans les caractères des espèces, dans les séquences d’ADN associées
%Pour reconstruire ces histoires, nous nous appuyons sur diverses sources de données telles que les caractères des espèces ou les séquences ADN associées à des objets présents ou fossiles.
%Le résultat est présenté sous la forme d’\emph{arbres phylogénétiques}, qui peuvent être comparés à des arbres généalogiques pour les objets étudiés.
%Les feuilles de ces arbres correspondent aux objets observés, et les autres nœuds à des ancêtres.

%Toutefois, contrairement aux études généalogiques, bien souvent en phylogénétique seules des informations sur des contemporains peuvent être obtenues.
%Pour reconstruire des arbres, la technique consiste à fixer une façon d’évaluer la distance qui sépare deux éléments.
%Une hypothèse de \emph{parcimonie} est ensuite faite, selon laquelle les éléments les moins distants sont supposés partager un ancêtre plus récent.
%Le défi restant consiste alors à trouver un arbre qui reflète cette hypothèse, citons à cet effet les algorithmes de \emph{neighbor joining}, UPGMA, …~\cite{Semple2003Phylogenetics}.

%Pour établir les liens entre les espèces, la phylogénétique s’est initialement fondée sur l’observation manuelle de leurs différences et similitudes morphologiques.
%Depuis la seconde moitié du XX\textsuperscript{e}~siècle, elle bénéficie de l’apport de nouvelles sources de données — les séquences protéiques et géniques — et de la contribution des sciences mathématiques, statistiques et informatiques pour le traiter ces données.
